\chapter{AI工具箱：理解你的协作伙伴}

本章旨在建立对 AI 工具的准确认知。AI 可以提高执行效率，但不能替代问题定义、领域判断与结果验证。

\section{大语言模型能做什么}

大语言模型可以协助完成代码生成、数据分析、结果解释和文献综述等任务。它尤其擅长根据明确要求快速生成候选方案，并帮助使用者探索不同的实现路径。

\section{大语言模型不能做什么}

\begin{enumerate}
  \item 它不会真正理解你的数据。
  \item 它不知道你的领域隐含假设。
  \item 它可能自信地犯错，即产生幻觉。
\end{enumerate}

例如，AI 不知道数据的采集机制、字段背后的业务含义，也不知道导师或审稿人可能质疑什么。AI 看到速度为 0 时，可能直接将其作为异常值删除；但在拥堵、排队和信号灯等待场景中，速度为 0 可能恰恰是重要信号。

\section{AI工具的生态全景：如何选择合适的AI工具}

\begin{enumerate}
  \item 对话型 AI，如 ChatGPT 和 Claude，适合问答、分析、解释和方案讨论。
  \item 代码协作型 AI，如 Codex、Claude Code 和 Copilot，适合在编辑器中结合项目上下文进行开发。
  \item 数据分析专用工具，如 Code Interpreter 和 Julius AI，适合直接上传数据并进行快速探索性数据分析。
  \item 专业领域工具，如 Overleaf AI 和 Elicit，适合论文写作、文献检索等特定场景。
\end{enumerate}

\section{你是项目负责人，AI是执行工程师}

不要只向 AI 索取答案，而应让 AI 参与解决问题的过程。可以让 AI 帮助列出检查清单、生成多个备选方案、指出方案风险，或扮演审稿人批评已有方法。目标、约束和最终判断仍应由人负责。

\section{实战练习}

\begin{enumerate}
  \item 判断什么样的问题不适合使用 AI 工具回答，以及使用 AI 后为什么可能达不到预期效果。
  \item 判断一个具体问题更适合使用对话型 AI，还是代码协作型 AI。
\end{enumerate}
